\begin{lemma}[Deficit estimate, algebraic form]
Let $m\ge2$ be an integer, let $S=\{0,1,\dots,m\}$ and let
$x_j=x^{(m)}_j=\cos(j\pi/m)$ be the Chebyshev extremal nodes of order $m$. Let
$P\in\mathbb{R}[X]$ satisfy
\[
  \deg P\le m,\qquad \lvert P(x)\rvert\le1\ \text{ for all }x\in[-1,1],\qquad
  P(1)=1 .
\]
Then
\[
  2\bigl(1+P(x_1)\bigr)\ \le\ \bigl(m^2-P'(1)\bigr)\,(1-x_1),
\]
where $P'$ is the formal derivative of $P$.
\end{lemma}
\begin{proof}
For $i\in S$ put
\[
  c_i=\Biggl(\prod_{j\in S\setminus\{i\}}(x_i-x_j)\Biggr)^{-1}
    \sum_{\substack{t\subseteq S\setminus\{i\}\\ \lvert t\rvert=m-1}}
      \ \prod_{a\in t}(1-x_a).
\]

\emph{Step 1: the interpolation formula.} The index set $S$ has
$\lvert S\rvert=m+1$ elements, and by the background proposition on strict
monotonicity of the extremal nodes the map $i\mapsto x_i$ is injective on $S$.
Let $Q\in\mathbb{R}[X]$ with $\deg Q\le m$, i.e. $\deg Q<\lvert S\rvert$. The
background Lagrange formula for an iterated derivative, applied with this $S$,
the node map $i\mapsto x_i$, the polynomial $Q$, the order $k=1$ (legitimate
since $1<m+1$) and the evaluation point $z=1$, gives, using $1!=1$ and
$\lvert S\rvert-(k+1)=(m+1)-2=m-1$,
\begin{equation}\label{eq:dp-interp}
  Q'(1)=\sum_{i\in S}
    \frac{Q(x_i)}{\prod_{j\in S\setminus\{i\}}(x_i-x_j)}
    \sum_{\substack{t\subseteq S\setminus\{i\}\\ \lvert t\rvert=m-1}}
      \ \prod_{a\in t}(1-x_a)
  =\sum_{i\in S}Q(x_i)\,c_i .
\end{equation}

\emph{Step 2: applying it to $P$ and to $T_m$.} The hypothesis $\deg P\le m$
lets us use \eqref{eq:dp-interp} for $Q=P$. The $m$-th Chebyshev polynomial
$T_m$ has $\deg T_m=m$ by the background Chebyshev proposition, so
\eqref{eq:dp-interp} also applies to $Q=T_m$. By the background propositions
$T_m(x_i)=(-1)^i$ for $i\in S$ and $T_m'(1)=m^2$. Subtracting the two instances
of \eqref{eq:dp-interp},
\begin{equation}\label{eq:dp-diff}
  m^2-P'(1)=T_m'(1)-P'(1)
   =\sum_{i\in S}\bigl((-1)^i-P(x_i)\bigr)c_i
   =\sum_{i\in S}\bigl(1-(-1)^iP(x_i)\bigr)\bigl((-1)^ic_i\bigr),
\end{equation}
the last equality holding term by term because $\bigl((-1)^i\bigr)^2=1$:
\[
  \bigl(1-(-1)^iP(x_i)\bigr)\bigl((-1)^ic_i\bigr)
  =(-1)^ic_i-\bigl((-1)^i\bigr)^2P(x_i)c_i
  =\bigl((-1)^i-P(x_i)\bigr)c_i .
\]

\emph{Step 3: every summand is nonnegative.} Fix $i\in S$. By the background
proposition on the extremal nodes, $x_i\in[-1,1]$, so the hypothesis gives
$\lvert P(x_i)\rvert\le1$; since
$(-1)^iP(x_i)\le\lvert(-1)^iP(x_i)\rvert=\lvert P(x_i)\rvert\le1$, we get
\[
  1-(-1)^iP(x_i)\ \ge\ 0 .
\]
By \noderef{LagrangeDerivSigns}, applied with this $m$, this $i$ (note $i\le m$)
and $k=m-1$, we have $(-1)^ic_i\ge0$. Hence each summand of \eqref{eq:dp-diff}
is a product of two nonnegative reals and so is nonnegative.

\emph{Step 4: the $i=0$ summand vanishes.} Since $x_0=\cos0=1$ and $(-1)^0=1$,
the $i=0$ summand of \eqref{eq:dp-diff} equals
$\bigl(1-P(1)\bigr)c_0=(1-1)c_0=0$ by the hypothesis $P(1)=1$. (This is the
paper's observation that the $j=0$ term of the sum vanishes because
$(T_m-P)(1)=0$.)

\emph{Step 5: dropping all summands but $i=1$.} By Steps 3 and 4,
\eqref{eq:dp-diff} is a sum over $S$ of nonnegative terms, so it is at least any
one of them. Note $1\in S$ because $m\ge2\ge1$. Taking the summand $i=1$, and
using $(-1)^1=-1$,
\[
  m^2-P'(1)\ \ge\ \bigl(1-(-1)P(x_1)\bigr)\bigl((-1)c_1\bigr)
   =\bigl(1+P(x_1)\bigr)\,(-c_1).
\]
By \noderef{LagrangeDerivCoeffOne} we have $c_1=\dfrac{-2}{1-x_1}$, hence
$-c_1=\dfrac{2}{1-x_1}$ and
\begin{equation}\label{eq:dp-key}
  m^2-P'(1)\ \ge\ \frac{2\bigl(1+P(x_1)\bigr)}{1-x_1}.
\end{equation}

\emph{Step 6: clearing the denominator.} Since $m\ge2$ we have
$0<\pi/m\le\pi/2$, and by the background proposition on strict monotonicity of
the extremal nodes $x_1<x_0=1$; hence $1-x_1>0$. Multiplying
\eqref{eq:dp-key} by the positive number $1-x_1$ preserves the inequality and
gives
\[
  \bigl(m^2-P'(1)\bigr)(1-x_1)\ \ge\ 2\bigl(1+P(x_1)\bigr),
\]
which is the assertion.
\end{proof}
